\subsection{Daftar Fitur yang Diimplementasikan}

Berikut adalah daftar fitur utama ResQfood beserta status implementasinya:

\begin{table}[H]
\centering
\caption{Status Implementasi Fitur}
\label{tab:implementasi}
\small
\begin{tabularx}{\textwidth}{lXX}
\toprule
\textbf{Modul} & \textbf{Fitur} & \textbf{Status} \\
\midrule
\multirow{4}{*}{Authentication} & Registrasi & Completed \\
& Login/Logout & Completed \\
& Forgot Password & Completed \\
& Role-based Access & Completed \\
\addlinespace
\multirow{4}{*}{Admin} & Dashboard Statistik & Completed \\
& Manajemen User & Completed \\
& Monitoring Produk & Completed \\
& Verifikasi Donasi & Partial \\
\addlinespace
\multirow{5}{*}{UMKM} & Kelola Produk & Completed \\
& Food Eligibility Assessment & Completed \\
& Food Eligibility Score & Completed \\
& Dynamic Pricing & Completed \\
& Riwayat Penjualan & Completed \\
\addlinespace
\multirow{4}{*}{Customer} & Marketplace & Completed \\
& Checkout \& Payment & Completed \\
& Tracking Pesanan & Completed \\
& Rating \& Ulasan & Completed \\
\bottomrule
\end{tabularx}
\end{table}

\subsection{Screenshot Antarmuka}

\subsubsection{Halaman Marketplace (Customer)}
\begin{figure}[H]
\centering
\fbox{\parbox{0.8\textwidth}{\centering\vspace{3cm}Screenshot Marketplace\\(Akan diganti dengan screenshot aktual)\vspace{3cm}}}
\caption{Tampilan Marketplace untuk Customer}
\label{fig:marketplace}
\end{figure}

Halaman marketplace menampilkan daftar makanan surplus dari UMKM terdekat dengan pengguna. Setiap produk menampilkan informasi harga asli, harga diskon, sisa waktu kelayakan, dan \textit{Food Trust Index}.

\subsubsection{Halaman Food Eligibility Assessment (UMKM)}
\begin{figure}[H]
\centering
\fbox{\parbox{0.8\textwidth}{\centering\vspace{3cm}Screenshot Food Eligibility\\(Akan diganti dengan screenshot aktual)\vspace{3cm}}}
\caption{Tampilan Food Eligibility Assessment}
\label{fig:food_eligibility}
\end{figure}

UMKM mengisi form penilaian kelayakan makanan yang mencakup waktu produksi, estimasi masa simpan, kondisi kemasan, dan suhu penyimpanan. Sistem kemudian menghitung \textit{Food Eligibility Score} secara otomatis.

\subsubsection{Halaman Dashboard Admin}
\begin{figure}[H]
\centering
\fbox{\parbox{0.8\textwidth}{\centering\vspace{3cm}Screenshot Dashboard Admin\\(Akan diganti dengan screenshot aktual)\vspace{3cm}}}
\caption{Tampilan Dashboard Admin}
\label{fig:dashboard_admin}
\end{figure}

Dashboard admin menampilkan ringkasan aktivitas platform meliputi jumlah UMKM aktif, total transaksi, volume \textit{food waste} yang berhasil diselamatkan, dan statistik donasi.

\subsection{Penjelasan Teknis Fitur Utama}

\subsubsection{Food Eligibility Assessment}
Sistem penilaian kelayakan makanan menggunakan model \textit{weighted scoring} dengan formula:

\begin{equation}
    FES = w_1 \cdot T_{prod} + w_2 \cdot S_{simpan} + w_3 \cdot K_{kemasan} + w_4 \cdot T_{suhu}
    \label{eq:fes}
\end{equation}

Dimana:
\begin{itemize}[leftmargin=*]
    \item $FES$ = Food Eligibility Score (0-100)
    \item $T_{prod}$ = Skor berdasarkan waktu produksi
    \item $S_{simpan}$ = Skor berdasarkan sisa masa simpan
    \item $K_{kemasan}$ = Skor berdasarkan kondisi kemasan
    \item $T_{suhu}$ = Skor berdasarkan suhu penyimpanan
    \item $w_1, w_2, w_3, w_4$ = Bobot masing-masing komponen
\end{itemize}

\subsubsection{Dynamic Pricing Engine}
Algoritma penetapan harga dinamis mempertimbangkan:

\begin{equation}
    P_{diskon} = P_{asal} \cdot \left(1 - \alpha \cdot \frac{T_{max} - T_{sisa}}{T_{max}}\right) \cdot \beta_{permintaan}
    \label{eq:dynamic_pricing}
\end{equation}

Dimana:
\begin{itemize}[leftmargin=*]
    \item $P_{diskon}$ = Harga diskon yang ditampilkan
    \item $P_{asal}$ = Harga asli produk
    \item $\alpha$ = Faktor diskon maksimal (0.5 = 50\%)
    \item $T_{max}$ = Masa simpan maksimal
    \item $T_{sisa}$ = Sisa waktu kelayakan
    \item $\beta_{permintaan}$ = Faktor permintaan lokal
\end{itemize}

\subsection{Integrasi Sistem}

\begin{itemize}[leftmargin=*]
    \item \textbf{Frontend-Backend} — Komunikasi via RESTful API dengan format JSON
    \item \textbf{Backend-Database} — Menggunakan Prisma ORM untuk akses PostgreSQL
    \item \textbf{Backend-ML Service} — HTTP request ke Python microservice untuk prediksi
    \item \textbf{Real-time} — Socket.IO untuk notifikasi dan update status pesanan
\end{itemize}
